

\section{Picard Iteration Design and Numerical Strategies}
In this section, we introduce the iteration design and numerical strategies.
The Picard hierarchy below is a numerical adaptation of the approximation
scheme in~\cite{wu-26}.  We preserve its principal dependency order---freezing
nonlinear coefficients at the previous iterate and successively solving
transport equations in their geometrically distinguished null directions---but
modify individual equations when an equivalent form is more suitable for
numerical integration.  Accordingly, the formulas in this section define the
numerical iteration used in the present work; they should not be read as a
literal transcription of the analytic scheme in~\cite{wu-26}.

\subsection{Picard iteration for Einstein equations}
In this section, we first consider Einstein scalar-field equations
\[\Rc=d\phi\otimes d\phi,\ \square_g\phi =0.\]

We work in region 
$$Q=\{-1<u<u_0,\ 0<v<v_0(u)\}.$$ 
The first task is to construct a valid initial data for ESE.
We directly prescribe $g,b,\Omega$ along $\Hb_0$. The only requirement is 
\[-\Omega\nabla_3(\Omega\tr\chib)-\frac{1}{2}(\Omega\tr\chib)^2-\left|\Omega\chibh\right|^2-4\Omega\omegab\Omega\tr\chib\geq 0\]
Then $\Omega e_3\phi(u,0)$ can be assigned via equation
\[\Omega e_3\phi=\left(-\Omega\nabla_3(\Omega\tr\chib)-\frac{1}{2}(\Omega\tr\chib)^2-\left|\Omega\chibh\right|^2-4\Omega\omegab\Omega\tr\chib\right)^{1/2}.\]
After defining $\phi(-1,0)$ and $\zeta(-1,0)$, we derive $\phi(u,0)$ from $\Omega e_3\phi(u,0)$ and $\zeta(u,0)$ via
\[\Omega\nabla_3\zeta+\frac{3}{2}\Omega\tr\chib\zeta+\Omega\chibh\cdot\zeta=-2\nabla(\Omega\omegab)-\dv(\Omega\chibh)+\frac{1}{2}\nabla(\Omega\tr\chib)-\Omega\tr\chib\nabla\log\Omega+\Omega e_3\phi\nabla\phi.\]
Thus $\eta=\zeta+\nabla\log\Omega$ and $\etab=-\zeta+\nabla\log\Omega$ are also prescribed.
We then define $\Omega\tr\chi(-1,0)$ and $\Omega\chih(-1,0)$, and compute $\Omega\chi(u,0)$ by 
\[\Omega\nabla_3(\Omega\tr\chi)+\Omega\tr\chib\Omega\tr\chi=\Omega^2\left(2\dv\eta+2\left|\eta\right|^2-2K(g)+|\nabla\phi|^2\right),\]
\[\Omega\nabla_3(\Omega\chih)+\frac{1}{2}\Omega\tr\chib\,\Omega\chih=\Omega^2\left(\nabla\hat\otimes\eta+\eta\hat\otimes\eta\right)-\frac{1}{2}\Omega\tr\chi\,\Omega\chibh+\frac{1}{2}\Omega^2\nabla\phi\hat\otimes\nabla\phi.\]
We also directly assign $\Omega e_4\phi(-1,0)$ and $\Omega\omega(-1,0)$, and compute them along $\Hb_0$ via
\[\Omega e_3(\Omega e_4\phi)+\frac{1}{2}\Omega\tr\chi\Omega e_3\phi+\frac{1}{2}\Omega\tr\chib\Omega e_4\phi-\Omega^2\Delta\phi-2\Omega^2\eta\cdot\nabla\phi=0,\]
\[\frac{1}{2}\Omega e_3\phi\Omega e_4\phi+\frac{1}{2}\Omega^2\left|\nabla\phi\right|^2=2\Omega\nabla_3(\Omega\omega)+\Omega^2K-\frac{1}{2}\Omega\chih\cdot\Omega\chibh+\frac{1}{4}\Omega\tr\chi\Omega\tr\chib-\Omega^2|\etab|^2+2\Omega^2\eta\cdot\etab.\]
Along $H_{-1}$, we prescribe $\Omega e_4\phi(-1,v)$ and $\Omega\omega(-1,v)$ to compute $\phi,\Omega$. Consider a symmetric and $g(-1,0)$-trace-free tensor $X_{AB}(-1,v)$.
We then prescribe $g,\Omega\chi$ along $H_{-1}$ by solving:
\begin{equation}
    \left\{\begin{aligned}
         & \Omega\chih_{AB}=\frac{1}{2}\left(X_{AC}\left(g(-1,0)\right)^{CD}{g}_{DB}+X_{BC}\left(g(-1,0)\right)^{CD}{g}_{DA}\right),\\
        &\partial_v (\Omega\tr\chi)=-\frac{1}{2}{\left(\Omega\tr\chi\right)}^2-4\Omega\omega\Omega\tr\chi-(\Omega e_4\phi)^2-\Omega\chih_{AB}\Omega\chih_{CD}g^{AC}g^{BD},\\
        &\Lie_v g_{AB}=\Omega\tr\chi g_{AB}+2\Omega\chih_{AB}.
    \end{aligned}\right.
\end{equation}

With these initial data, we can start the iteration.
The quantities to be iterated are 
\[g^{(i)},(b^{(i)})^A,\log\Omega^{(i)},(\Omega\chi)^{(i)}_{AB},(\Omega\chib)^{(i)}_{AB},(\zeta^{(i)})^A,(\Omega\omega)^{(i)},(\Omega\omegab)^{(i)},(\Omega e_3\phi)^{(i)}, (\Omega e_4\phi)^{(i)},\nabla_A\phi^{(i)} .\]

Initially, we set $g^{(0)}=g|_{v=0}$, $\Omega^{(0)}=\Omega|_{v=0}$, $\zeta^{(0)}=\zeta|_{v=0}$,
\[
 (b^{(0)})^A(u,v)=b^A(u,0)-\int_0^v4(\Omega^{(0)})^2
 (\zeta^{(0)})^A(u,v')\,\dd v',
\]
$\chih^{(0)}=\chih|_{v=0}$, $\tr\chi^{(0)}=\tr\chi|_{v=0}$, $\chibh^{(0)}=\chibh|_{v=0}$, $\tr\chib^{(0)}=\tr\chib|_{v=0}$, $\omega^{(0)}=\omega|_{v=0}$, $\omegab^{(0)}=\omegab|_{v=0}$, $(e_3\phi)^{(0)}=e_3\phi|_{v=0}$, $(e_4\phi)^{(0)}=e_4\phi|_{v=0}$, $(\nabla_A\phi)^{(0)}=\nabla_A\phi|_{v=0}$.
For all $i$, we always require 
\[\eta^{(i)}_A=(\zeta^{(i)})^B g^{(i)}_{AB}+\nabla_A\log\Omega^{(i)},\ \etab^{(i)}_A=-(\zeta^{(i)})^B g^{(i)}_{AB}+\nabla_A\log\Omega^{(i)}\]

\noindent\textbf{1. Construction of $(\Omega\chih)^{(i+1/2)}$, $\partial_v\phi^{(i+1)}$, and $\nabla_A\phi^{(i+1)}$}

We use \eqref{eq:Ricci_coefficients_propagation_2} to define $(\Omega\chih)^{(i+1/2)}$,
\begin{equation}
    \left\{
    \begin{aligned}
        & \Lie_{\partial_u+b^{(i)}}(\Omega\chih^{(i+1/2)})_{AB}-\frac{1}{2}(\Omega\tr\chib)^{(i)}(\Omega\chih^{(i+1/2)})_{AB}\\
        &\qquad\qquad -(\Omega\chibh^{(i)})_A^{\quad C}(\Omega\chih^{(i+1/2)})_{CB}-(\Omega\chibh^{(i)})_B^{\quad C}(\Omega\chih^{(i+1/2)})_{CA} \\
        &\qquad =(\Omega^{(i)})^2\left(\left(\nabla^{(i)}{\hat\otimes}\eta^{(i)}+\eta^{(i)}{\hat\otimes}\eta^{(i)}+\frac{1}{2}\nabla\phi^{(i)}\hat\otimes\nabla\phi^{(i)}\right)_{AB}-\frac{1}{2}(\tr\chi)^{(i)}{\chibh^{(i)}}_{AB}\right),\\
        &{\Omega\chih^{(i+1/2)}}_{AB}(-1,v)=\frac{1}{2}\left(\Omega\chih_{AC}g^{CD}{g^{(i)}}_{DB}+\Omega\chih_{BC}g^{CD}{g^{(i)}}_{DA}\right)(-1,v).
    \end{aligned}
    \right.
\end{equation}
Using \eqref{eq:ese-wave-eq}, we define $\partial_v\phi^{(i+1)}$ by
\begin{equation}
    \left\{\begin{aligned}
        &(\Omega e_3)^{(i)}(\partial_v\phi^{(i+1)})+\frac{1}{2}(\Omega\tr\chib)^{(i)}(\partial_v\phi^{(i+1)})\\
        &\qquad =(\Omega^{(i)})^2\Delta^{(i)}\phi^{(i)}-\frac{1}{2}(\Omega\tr\chi)^{(i)}(\Omega e_3)^{(i)}\phi^{(i)}+2(\Omega^{(i)})^2\eta^{(i)}\nabla\phi^{(i)},\\
        &\partial_v\phi^{(i+1)}(-1,v)=\partial_v\phi(-1,v).
    \end{aligned}\right.
\end{equation}
It follows that
\[\nabla_A\phi^{(i+1)}(u,v)=\nabla_A\phi(u,0)+\partial_A\left(\int_0^v\partial_v\phi^{(i+1)}\right).\]


\noindent\textbf{2. Construction of $(\Omega\omegab)^{(i+1/2)}$}

We use \eqref{eq:Ricci_coefficients_propagation_3} to define $(\Omega\omegab)^{(i+1/2)}$,
\begin{equation}
    \left\{
    \begin{aligned}
        &4\partial_v(\Omega\omegab)^{(i+1/2)}=\Omega\chih^{(i+1/2)}_{AB}\Omega\chibh^{(i)}_{CD}(g^{(i)})^{AC}(g^{(i)})^{BD}+\frac{1}{2}(\Omega\tr\chi)^{(i)}(\Omega\tr\chib)^{(i)}\\
        &\qquad-4(\Omega^{(i)})^2\eta^{(i)}\cdot\etab^{(i)}+(\partial_u+b^{(i)})(\Omega\tr\chi)^{(i)}-2(\Omega^{(i)})^2\dv^{(i)}\eta^{(i)}\\
        &\qquad+(\partial_v\phi)^{(i+1)}(\Omega e_3\phi)^{(i)},\\
        &(\Omega\omegab)^{(i+1/2)}(u,0)=\Omega\omegab(u,0).
    \end{aligned}\right.
\end{equation}

\noindent\textbf{3. Construction of $\Omega^{(i+1)}$, $\nabla_A\log\Omega^{(i+1)}$, and $(\Omega\omega)^{(i+1)}$}

With $(\Omega\omegab)^{(i+1/2)}$, we can integrate it to obtain $\Omega^{(i+1)}$
\begin{equation} 
    \left\{
    \begin{aligned}
        &\left(\partial_u+b^{(i)}\right)\log \Omega^{(i+1)}=-2(\Omega\omegab)^{(i+1/2)},\\
        &\log\Omega^{(i+1)}(-1,v)=\log\Omega(-1,v),
    \end{aligned}
    \right.
\end{equation}
\begin{equation}
    \nabla_A\log\Omega^{(i+1)}=\partial_A \log\Omega^{(i+1)}.
\end{equation}
We construct $(\Omega\omega)^{(i+1)}$ by integration instead of differentiation, 
\begin{equation} 
    \left\{\begin{aligned}
        4\left(\partial_u+b^{(i)}\right)&(\Omega\omega)^{(i+1)}=\Omega\chih^{(i+1/2)}_{AB}\Omega\chibh^{(i)}_{CD}(g^{(i)})^{AC}(g^{(i)})^{BD}+\frac{1}{2}(\Omega\tr\chi)^{(i)}(\Omega\tr\chib)^{(i)}\\
        &-4(\Omega^{(i)})^2\eta^{(i)}\cdot\etab^{(i)}+(\partial_u+b^{(i)})(\Omega\tr\chi)^{(i)}-2(\Omega^{(i)})^2\dv^{(i)}\eta^{(i)}\\
        &+(\partial_v\phi)^{(i+1)}(\Omega e_3\phi)^{(i)}-8(\Omega^{(i)})^2(\zeta^{(i)})^A\nabla_A\log\Omega^{(i+1)},\\
        (\Omega\omega)^{(i+1)}&(-1,v)=\Omega\omega(-1,v).
    \end{aligned}\right.
\end{equation}
We note that the construction of $(\Omega\omega)^{(i+1)}$ is equivalent to $(\Omega\omega)^{(i+1)}=-\frac{1}{2}\partial_v\log\Omega^{(i+1)}$, because 
\begin{equation}
   \begin{aligned}
     4(\Omega e_3)^{(i)}(\Omega\omega)^{(i+1)}=&-2\partial_v(\Omega e_3)^{(i)}\log\Omega^{(i+1)}+2\Lie_v b^{(i)}\nabla\log\Omega^{(i+1)}\\
     =&4\partial_v(\Omega\omegab)^{(i+1/2)}-8(\Omega^{(i)})^2(\zeta^{(i)})^A\nabla_A\log\Omega^{(i+1)}.
   \end{aligned}
\end{equation}

\noindent\textbf{4. Construction of $\zeta^{(i+1)}$, $b^{(i+1)}$, and $(\Omega\omegab)^{(i+1)}$}

From \eqref{eq:Ricci_coefficients_propagation_4}, we define $\zeta^{(i+1)}$,
\begin{equation} 
    \left\{\begin{aligned}
        &\Lie_v(\zeta^{(i+1)})^A=-2(\Omega\tr\chi)^{(i)}(\zeta^{(i)})^A-2(g^{(i)})^{AC}\left((\Omega\chih)^{(i+1/2)}\right)_{CB}(\zeta^{(i)})^B\\
        &\qquad+2(g^{(i)})^{AB}\partial_B(\Omega\omega)^{(i+1)} +(\dv)^{(i)}(\Omega\chih)^{(i+1/2)}_B(g^{(i)})^{BA} \\
        &\qquad-\frac{1}{2}(g^{(i)})^{AB}\partial_B(\Omega\tr\chi)^{(i)}+(\Omega\tr\chi)^{(i)}(g^{(i)})^{AB}\partial_B\log\Omega^{(i+1)}\\
        &\qquad-(\partial_v\phi)^{(i+1)}\partial_B\phi^{(i)}(g^{(i)})^{AB},\\
        &(\zeta^{(i+1)})^A(u,0)=\zeta^A(u,0).
    \end{aligned}\right.
\end{equation}
We remark that it is necessary to treat $\zeta^{(i+1)}$ as a vector field instead of a 1-form in order to define $b^{(i+1)}$,
\begin{equation} 
    \left\{\begin{aligned}
        &\Lie_v (b^{(i+1)})^A=-4(\Omega^{(i+1)})^2(\zeta^{(i+1)})^A,\\
        &(b^{(i+1)})^A(u,0)=b^A(u,0).
    \end{aligned}\right.
\end{equation}

The incoming lapse coefficient must then use the updated shift:
\begin{equation}
    (\Omega\omegab)^{(i+1)}=(\Omega\omegab)^{(i+1/2)}
    -\frac12\bigl(b^{(i+1)}-b^{(i)}\bigr)\cdot\nabla\log\Omega^{(i+1)}.
\end{equation}
The numerical implementation applies the angular projection to this correction
and restores the prescribed traces after projection and relaxation.

\noindent\textbf{5. Construction of $g^{(i+1)}$ and $(\Omega\chi)^{(i+1)}$}

It remains to define $g^{(i+1)}$. We consider the following $(g,\tr\chi,\chih)$ system,
\begin{equation}
    \left\{
    \begin{aligned}
        &\Lie_v(g^{(i+1)})_{AB}=(\Omega\tr\chi)^{(i+1)}(g^{(i+1)})_{AB}+2{(\Omega\chih)^{(i+1)}}_{AB},\\
        &\partial_v\left(\Omega\tr\chi\right)^{(i+1)}=-\frac{1}{2}\left((\Omega\tr\chi)^{(i+1)}\right)^2-4(\Omega\omega)^{(i+1)}(\Omega\tr\chi)^{(i+1)}\\
        &\qquad\qquad- {(\Omega\chih)^{(i+1)}}_{AB}{(\Omega\chih)^{(i+1)}}_{CD}{g^{(i+1)}}^{AC}{g^{(i+1)}}^{BD}-(\partial_v\phi^{(i+1)})^2,\\
        &(\Omega\chih)^{(i+1)}_{AB}=\frac{1}{2}\left((\Omega\chih)^{(i+1/2)}_{AC}{(g^{(i)})}^{CD}{g^{(i+1)}}_{DB}
        +(\Omega\chih)^{(i+1/2)}_{BC}{(g^{(i)})}^{CD}{g^{(i+1)}}_{DA}\right),\\
        &(g^{(i+1)})_{AB}(u,0)=g_{AB}(u,0),\ \ (\Omega\tr\chi)^{(i+1)}(u,0)=\Omega\tr\chi(u,0).
    \end{aligned}
    \right.
\end{equation}

\noindent\textbf{6. Construction of $(\Omega\chib)^{(i+1)}$ and $(\Omega e_3\phi)^{(i+1)}$}

There are two equivalent ways to define $(\Omega\chib)^{(i+1)}$. The first one is using $\Omega\chib_{AB}=\frac{1}{2}\Lie_{\Omega e_3}g_{AB}$. We use the second way, which is using equation
\begin{equation}
    \left\{\begin{aligned}
        &\Lie_{v}(\Omega\chib)^{(i+1)}_{AB}= \Lie_{\partial_u+b^{(i+1)}}(\Omega\chi)^{(i+1)}_{AB}-2(\Omega^{(i+1)})^2\left(\nabla^{(i+1)}_A\zeta^{(i+1)}_B+\nabla^{(i+1)}_B\zeta^{(i+1)}_A\right)\\
        &\qquad\qquad-4(\Omega^{(i+1)})^2\left(\zeta^{(i+1)}_A\nabla_B\log\Omega^{(i+1)}+\zeta^{(i+1)}_B\nabla_A\log\Omega^{(i+1)}\right),\\
        &(\Omega\chib)^{(i+1)}_{AB}(u,0)= \Omega\chib_{AB}(u,0).
    \end{aligned}\right.
\end{equation}
We remark that
\begin{equation}
    \begin{aligned}
        &\Omega\nabla_4(\Omega\chib)_{AB}-\Omega\nabla_3(\Omega\chi)_{AB}\\
        =&\Lie_{v}(\Omega\chib)_{AB}-(\Omega\chi)_{A}^{\ C}(\Omega\chib)_{CB}-(\Omega\chi)_{B}^{\ C}(\Omega\chib)_{CA}\\
        &-\Lie_{\partial_u+b}(\Omega\chi)_{AB}+(\Omega\chib)_{A}^{\ C}(\Omega\chi)_{CB}+(\Omega\chib)_{B}^{\ C}(\Omega\chi)_{AC}\\
        =&\Lie_{v}(\Omega\chib)_{AB}-\Lie_{\partial_u+b}(\Omega\chi)_{AB}.
    \end{aligned}
\end{equation}
Theoretically, $(\Omega\chib)_{AB}^{(i+1)}=\frac{1}{2}\Lie_{\partial_u+b^{(i+1)}}g^{(i+1)}_{AB}.$ Similarly, we construct $(\Omega e_3\phi)^{(i+1)}$ by
\begin{equation}
    \left\{\begin{aligned}
        &\partial_v(\Omega e_3\phi)^{(i+1)}=(\partial_u+b^{(i+1)})\partial_v\phi^{(i+1)}-4(\Omega^{(i+1)})^2(\zeta^{(i+1)})^A\nabla_A\phi^{(i+1)},\\
        &(\Omega e_3\phi)^{(i+1)}(u,0)=\Omega e_3\phi(u,0).
    \end{aligned}\right.
\end{equation}

For Einstein vacuum equation
\[\Rc_{\mu\nu}=0,\]
removing all $\phi$-terms of ESE gives the iterative construction directly, so we omit the repeation of construction. 


\subsection{Curvature residues}
Computing curvature usually involves second order derivatives, which would heavily reduce the numerical precision. 
Especially, $\lim_{v\rightarrow 0}R_{4A4B}(-1,v)$ can be infinite for short pulse perturbation and $\Rc_{44}$ is hard to compute in this way, even though it is theoretically zero in our construction.
See \cite{an-25} for the example of short pulse with infinite $R_{4A4B}$.
We use intrinsic equations of Lorentzian manifold to compute the Ricci curvature:
\[\Omega^2\Rc_{33}=-\Omega\nabla_3(\Omega\tr\chib)-4\Omega\omegab\Omega\tr\chib-\frac{1}{2}(\Omega\tr\chib)^2-|\Omega\chibh|^2,\]
\[\Omega\Rc_{3A}=\Omega\nabla_3\zeta+\frac{3}{2}\Omega\tr\chib\zeta+\Omega\chibh\cdot\zeta+2\nabla(\Omega\omegab)+\dv(\Omega\chibh)-\frac{1}{2}\nabla(\Omega\tr\chib)+\Omega\tr\chib\nabla\log\Omega,\]
\[\Omega\Rc_{4A}=-\Omega\nabla_4\zeta-\frac{3}{2}\Omega\tr\chi\zeta-\Omega\chih\cdot\zeta+2\nabla(\Omega\omega)+\dv(\Omega\chih)-\frac{1}{2}\nabla(\Omega\tr\chi)+\Omega\tr\chi\nabla\log\Omega,\]
\[\Omega^2\widehat{\Rc}_{AB}=\Omega\nabla_3(\Omega\chih)+\frac{1}{2}\Omega\tr\chib\,\Omega\chih-\Omega^2\left(\nabla\hat\otimes\eta+\eta\hat\otimes\eta\right)+\frac{1}{2}\Omega\tr\chi\,\Omega\chibh,\]
\[ \Omega^2g^{AB}\Rc_{AB}=\Omega\nabla_3(\Omega\tr\chi)+\Omega\tr\chi\Omega\tr\chib-2\Omega^2\dv\eta-2\Omega^2|\eta|^2+2\Omega^2K,\]
\[\Omega^2\Rc_{34}=4\Omega\nabla_4(\Omega\omegab)-\Omega\chih\cdot\Omega\chibh-\frac{1}{2}\Omega\tr\chi\Omega\tr\chib+4\Omega^2\eta\cdot\etab -\Omega\nabla_3(\Omega\tr\chi)+2\Omega^2\dv\eta.\]
For a fresh check of the outgoing constraint we also use
\[
 \Omega^2\Rc_{44}=-\partial_v(\Omega\tr\chi)
 -4(\Omega\omega)(\Omega\tr\chi)
 -\tfrac12(\Omega\tr\chi)^2-|\Omega\chih|^2.
\]
The coefficient $4$ includes the product-rule contribution from weighting
$\tr\chi$ by $\Omega$, since $\partial_v\log\Omega=-2\Omega\omega$.
Thus no outgoing null derivative of $\chih$ is required.  These formulas use
first null derivatives of the stored connection variables; the intrinsic
Gauss curvature $K$ uses the angularly smooth section metric.  The numerical
metric--connection relations are checked separately.  Endpoint values at
singular coordinate Jacobians are excluded from the reported residuals.

\subsection{Legendre--Gauss--Lobatto spectral elements}

In the experiments, we will add non-spherical symmetric perturbation to the Einstein scalar-field or vacuum systems. To make the metric less regular, the perturbation is required to be of size $O(v^\delta)$.
To deal with the non-smoothness, the characteristic coordinates are first replaced by coordinates in which the
corner profiles are smoother.  In particular, we use
\[
    \tau=-\log(-u),\qquad
    s=\left(\frac{v}{v_{\max}}\right)^\delta ,
    \qquad
    u=-e^{-\tau},\qquad
    v=v_{\max}s^{1/\delta}.
\]
Thus a profile proportional to $v^\delta$ is linear in $s$.  This change of
variables does not regularize the physical solution; it only represents its
fractional behavior by a smooth function of the computational coordinate.
Physical derivatives and integrals are recovered by
\[
    \partial_u=\frac{1}{-u}\partial_\tau,\qquad
    \partial_v=\left(\frac{dv}{ds}\right)^{-1}\partial_s,\qquad
    \int_0^v F(v')\,dv'
    =\int_0^s F(v(s'))\frac{dv}{ds'}\,ds'.
\]

We divide each computational coordinate interval into elements.  On an
element of degree $p$, let $\{\xi_j\}_{j=0}^p$ be the
Legendre--Gauss--Lobatto nodes, characterized by
\[
    (1-\xi_j^2)P_p'(\xi_j)=0,
\]
including the two endpoints.  A component $F$ is represented by its Lagrange
interpolant
\[
    I_pF(\xi)=\sum_{j=0}^p F(\xi_j)\ell_j(\xi).
\]
The element differentiation and indefinite-integration matrices are
\[
    D_{ij}=\ell_j'(\xi_i),\qquad
    Q_{ij}=\int_{-1}^{\xi_i}\ell_j(\xi)\,d\xi.
\]
Consequently, interpolation of intermediate Runge--Kutta stages,
differentiation, and cumulative integration are all obtained from the same
element polynomial.  Neighboring elements share their endpoint value, while
the two derivative traces at an interior interface are combined into a
single interface derivative.  This composite construction permits local
refinement in either characteristic direction without replacing the
underlying transport hierarchy.

\subsection{Tensor spherical-harmonic Galerkin discretization}
The angular variables are represented on a quasi-uniform, pole-free point set
on $\mathbb S^2$.  Geometric vectors and tensors are stored as ambient
Cartesian tensors and projected to the tangent bundle.  This avoids the
coordinate singularities of a polar chart.  For a scalar field, we use the
real spherical-harmonic expansion
\[
    f_L(\vartheta)
    =\sum_{\ell=0}^{L}\sum_{m=-\ell}^{\ell}
      f_{\ell m}Y_{\ell m}(\vartheta).
\]
The coefficients are obtained by an overdetermined least-squares analysis of
the nodal values, and angular derivatives are computed by differentiating the
harmonic basis.

Vectors and symmetric two-tensors are not expanded componentwise as
unrelated scalars.  We use the electric and magnetic vector harmonics
\[
    Y_A^{E,\ell m}\sim\nabla_A Y_{\ell m},\qquad
    Y_A^{B,\ell m}\sim\epsilon_A{}^B\nabla_B Y_{\ell m},
\]
and the trace, electric, and magnetic tensor harmonics
\[
    \gamma_{AB}Y_{\ell m},\qquad
    Y_{AB}^{E,\ell m}\sim
       \left(\nabla_A\nabla_B Y_{\ell m}\right)^{\mathrm{TF}},
    \qquad
    Y_{AB}^{B,\ell m}\sim
       \epsilon_{(A}{}^C\nabla_{B)}\nabla_CY_{\ell m}.
\]
These typed spaces preserve tangency, symmetry, and the tensorial
transformation law.  Trace-free quantities are projected with respect to the
current section metric,
\[
    T_{AB}^{\mathrm{TF}}
    =T_{AB}-\frac12(\tr_\gamma T)\gamma_{AB},
\]
rather than with respect to a fixed background metric.

Let $\mathcal V_L$ denote the retained harmonic space and let
$\mathcal V_W$, with $W>L$, be a work space used to form angular products.
For a semidiscrete unknown $U_L\in\mathcal V_L$, every nonlinear right-hand
side is defined by
\[
    \partial U_L=\Pi_L F(U_L),
\]
where $F(U_L)$ is evaluated in the work space and $\Pi_L$ is the appropriate
scalar, vector, or tensor Galerkin projection.  The retained/work separation
reduces angular aliasing while keeping the evolved state in one declared
finite-dimensional geometric space.
